\documentclass{beamer}
\usetheme{Madrid}
\usecolortheme{dove}

\usepackage{amsmath,amssymb}
\usepackage{booktabs}
\usepackage{hyperref}

\title{CNN-Based Semiconductor Wafer Defect Detection}
\subtitle{WM-811K Dataset Classification Project}
\author{Anindita Paul \and Brandon Pyle \and Anand Rajan \and Brett Rettura}
\institute{AI 570 - Team 4}
\date{\today}

\begin{document}

\frame{\titlepage}

\begin{frame}
  \frametitle{Problem}
  Semiconductor wafer defect classification on the WM-811K dataset.

  \begin{itemize}
    \item $\sim$120K samples, 9 defect classes
    \item Severe class imbalance ($\sim$85\% ``none'')
    \item Variable wafer map shapes $\rightarrow$ common resize
  \end{itemize}
\end{frame}

\begin{frame}
  \frametitle{Approach}
  Compare three architectures on the same preprocessed dataset:

  \begin{itemize}
    \item Custom CNN (4-block, from scratch, $\sim$1.2M params)
    \item ResNet-18 transfer learning (ImageNet pretrained, $\sim$11.2M params)
    \item EfficientNet-B0 transfer learning ($\sim$5.3M params)
  \end{itemize}

  \vspace{0.3cm}
  Class-weighted cross-entropy for imbalance; stratified train/val/test split.
\end{frame}

\begin{frame}
  \frametitle{Repository Layout}
  \small
  \begin{center}
    \begin{tabular}{ll}
      \toprule
      Module & Role \\
      \midrule
      \texttt{src/data} & loading, preprocessing, transforms \\
      \texttt{src/models} & CNN, ResNet, EfficientNet, ViT, Swin \\
      \texttt{src/training} & trainer, config, DDP, SimCLR, SupCon \\
      \texttt{src/inference} & FastAPI server, Grad-CAM, uncertainty \\
      \texttt{src/analysis} & metrics, visualization, OOD, anomaly \\
      \texttt{src/federated} & federated averaging client/server \\
      \texttt{src/augmentation} & synthetic wafer-map generation \\
      \texttt{tests} & unit + integration \\
      \bottomrule
    \end{tabular}
  \end{center}
\end{frame}

\begin{frame}
  \frametitle{Preprocessing}
  \begin{itemize}
    \item Resize all wafer maps to a common resolution (default $96\times96$)
    \item Normalize per-channel
    \item ImageNet statistics for pretrained models only
    \item Optional augmentation: flip, rotate, affine
  \end{itemize}
\end{frame}

\begin{frame}
  \frametitle{Training}
  \begin{itemize}
    \item AdamW + cosine schedule
    \item Mixed precision on GPU
    \item Layer freezing for transfer-learning backbones
    \item Early stopping on validation macro-F1
  \end{itemize}

  \vspace{0.2cm}
  Entry point: \texttt{python train.py --model all --epochs 20 --device cuda}
\end{frame}

\begin{frame}
  \frametitle{Evaluation}
  Metrics reported per model:
  \begin{itemize}
    \item Accuracy and weighted F1 (dominated by majority class)
    \item Macro F1 (primary metric for imbalanced setting)
    \item Per-class precision / recall / F1
    \item Confusion matrix visualization
    \item Grad-CAM overlays for spatial interpretability
  \end{itemize}
\end{frame}

\begin{frame}
  \frametitle{Extensions}
  Beyond the three baselines, the repository supports:
  \begin{itemize}
    \item Monte Carlo dropout uncertainty estimation
    \item OOD / anomaly detection (IsolationForest, OC-SVM, autoencoder)
    \item Federated averaging simulation
    \item Synthetic data augmentation for minority classes
    \item Model ensembling (voting, averaging, weighted)
    \item SimCLR / SupCon self-supervised pretraining
    \item Distributed multi-GPU training (DDP)
  \end{itemize}
\end{frame}

\begin{frame}
  \frametitle{Running the Project}
  \small
  \begin{enumerate}
    \item Clone and install: \texttt{pip install -e ".[dev]"}
    \item Place dataset at \texttt{data/LSWMD\_new.pkl}
    \item Run tests: \texttt{pytest -q}
    \item Train: \texttt{python train.py --model all --epochs 20}
    \item Evaluate saved checkpoints or run Colab notebook
  \end{enumerate}
\end{frame}

\begin{frame}
  \frametitle{Known Limitations}
  \begin{itemize}
    \item Trained checkpoints not versioned in repo (regenerate via training)
    \item Inference server exposes a single active model at a time
    \item Final accuracy numbers depend on dataset split and epoch count
  \end{itemize}
\end{frame}

\begin{frame}
  \frametitle{Summary}
  \begin{itemize}
    \item Modular PyTorch pipeline for WM-811K defect classification
    \item Three baseline architectures with transfer-learning + scratch CNN
    \item Full evaluation + explainability (Grad-CAM) + uncertainty tooling
    \item Reproducible via \texttt{pytest} and a Colab-ready notebook
  \end{itemize}
\end{frame}

\end{document}
